\begin{center}
    \section*{Priedai}

\end{center}
\addcontentsline{toc}{section}{Priedai}

\newcounter{priedas}
\renewcommand{\thepriedas}{\arabic{priedas}}

\newcommand{\priedas}[2]{%
    \refstepcounter{priedas}%
    \par\noindent\subsection*{\thepriedas\ priedas. #1}%
    \label{#2}%
    \addcontentsline{toc}{subsection}{\thepriedas\ priedas. #1}%
}

\priedas{Kodo pavyzdys}{priedas:veiklos}

\begin{inline_code}
    for i in range(10):
    print(i)
\end{inline_code}


% \code{kodas/adaptive.lua}{luastyle}
% \code{kodas/example.py}{pythonstyle}
\code{kodas/example.cs}{csharpstyle}


\newpage


\priedas{Citavimas}{priedas:citavimas}

Bibliografijos įrašai saugomi \textit{bakis.bib} faile.

Apibrėžtus įrašus galima cituoti tekste naudojant \textit{cite} komandą, pavyzdžiui: \textit{\cite{lituanistika_gudinavicius_2021}} arba \cite{capitaineplanete_pr,all_matches_function}.

\priedas{Paveiklai}{priedas:uzduotis}

\ref{fig:uzduotis} paveiksle pateikiamas siekiamas programos išvesties rezultatas.

\begin{figure}[h]
    \centering
    \includegraphics[width=0.5\textwidth]{img/uzduotis.png}
    \caption{Užduotis}
    \label{fig:uzduotis}
\end{figure}

\priedas{Lentelės}{priedas:lentele}


\begin{table}[h]
    \centering
    \begin{tabular}{|c|c|}
        \hline
        Stulpelis 1 & Stulpelis 2 \\
        \hline
        Reikšmė 1   & Reikšmė 2   \\
        \hline
    \end{tabular}
    \caption{Lentelės pavyzdys}
    \label{tab:lentele}
\end{table}


\newpage